\chapter{Conception de la solution}

\section{Architecture générale}

L'architecture sépare la collecte industrielle, le transport, le traitement,
le stockage et la visualisation. Cette séparation facilite les essais et
permet de remplacer le simulateur par Litmus Edge sans réécrire les tableaux
de bord.

\begin{figure}[H]
  \centering
  \includegraphics[width=\textwidth]{architecture-industrielle.pdf}
  \caption{Architecture fonctionnelle de GlassInk Monitor}
  \label{fig:architecture}
\end{figure}

Le PLC et l'imprimante appartiennent à la couche de contrôle existante.
Litmus Edge fournit le point d'intégration. Mosquitto transporte les messages.
Node-RED vérifie leur structure, les convertit en Line Protocol et les écrit
dans InfluxDB. Grafana interroge ensuite le bucket de supervision.

\section{Choix de la pile technique}

\begin{center}
  \stackicon{python.pdf}{Python\\simulateur}
  \hfill
  \stackicon{eclipsemosquitto.pdf}{Mosquitto\\MQTT}
  \hfill
  \stackicon{nodered.pdf}{Node-RED\\ingestion}
  \hfill
  \stackicon{influxdb.pdf}{InfluxDB\\stockage}
  \hfill
  \stackicon{grafana.pdf}{Grafana\\supervision}
  \hfill
  \stackicon{docker.pdf}{Docker\\déploiement}
\end{center}

\subsection{MQTT et Mosquitto}

MQTT suit le modèle publication/abonnement. Le producteur publie sur un topic,
le courtier distribue le message et les consommateurs s'abonnent aux topics
utiles. Mosquitto a été retenu parce que le volume du projet reste faible et
qu'il fournit les fonctions nécessaires : authentification, ACL, TLS,
persistance et outils de test~\cite{mosquitto}.

Le prototype utilise le QoS 1. Ce niveau garantit une livraison ``au moins une
fois''~\cite{mqtt5}. Il améliore la fiabilité, mais autorise un doublon lors
d'une retransmission. L'application doit donc être idempotente.

\subsection{Node-RED}

Node-RED rend le chemin des données visible dans un éditeur. Ce choix est
adapté au passage du simulateur vers Litmus Edge : le mapping peut être
inspecté et modifié sans cacher le traitement dans un service opaque. Les flux
restent versionnés dans \texttt{flows.json}. L'éditeur est authentifié et les
connexions utilisent les variables d'environnement~\cite{nodered}.

\subsection{InfluxDB}

La plateforme produit surtout des données horodatées : mesures continues,
changements d'état, défauts et résultats de marquage. InfluxDB 2.7 les conserve
dans le bucket \texttt{printer\_monitoring}. Trois mesures structurent le
contenu : \texttt{printer\_telemetry}, \texttt{printer\_event} et
\texttt{marking\_event}. Les identifiants utilisés pour filtrer deviennent des
tags ; les valeurs numériques et descriptives restent des champs
~\cite{influxschema}.

\subsection{Grafana}

Grafana interroge nativement InfluxDB. Il affiche les séries temporelles, les
tableaux et les indicateurs calculés avec Flux. Les règles d'alerte réutilisent
les mêmes données~\cite{grafanainflux}. Les tableaux de bord et
la source de données sont provisionnés par fichiers, ce qui évite une
configuration manuelle lors d'une nouvelle installation.

\subsection{Docker Compose}

Chaque composant tourne dans son conteneur. Le fichier Compose décrit les
services, leurs réseaux, leurs volumes, leurs variables et leurs contrôles de
santé. Cette approche reproduit le même environnement sur un poste de
développement et sur une VM de démonstration~\cite{dockercompose}.

\section{Cycle de vie des données}

\begin{figure}[H]
  \centering
  \includegraphics[width=\textwidth]{pipeline-donnees.pdf}
  \caption{Traitement d'une mesure entre la machine et le tableau de bord}
  \label{fig:pipeline}
\end{figure}

Chaque message reçoit un \texttt{message\_id}. Node-RED contrôle les champs
obligatoires, construit un point InfluxDB puis l'envoie à l'API d'écriture v2.
La clé d'un point associe la mesure, son jeu de tags et son horodatage. Une
retransmission strictement identique réécrit donc le même point au lieu d'en
créer un second~\cite{influxduplicate}. Cette propriété couvre les doublons
attendus avec le QoS~1. L'écriture de plusieurs mesures n'est toutefois pas une
transaction globale ; en cas d'échec partiel, Node-RED rejoue les points avec
les mêmes clés.

\section{Topics MQTT}

\begin{table}[H]
\centering
\caption{Arborescence MQTT du prototype}
\begin{tabularx}{\textwidth}{p{0.27\textwidth}Y}
\toprule
Type & Topic \\
\midrule
Télémétrie &
\texttt{sgx/<site>/<ligne>/printer/<id>/telemetry} \\
Événement &
\texttt{sgx/<site>/<ligne>/printer/<id>/event} \\
Marquage &
\texttt{sgx/<site>/<ligne>/printer/<id>/marking} \\
Rejet &
\texttt{sgx/system/dead-letter} \\
État Node-RED &
\texttt{sgx/system/node-red/status} \\
\bottomrule
\end{tabularx}
\end{table}

Lors du raccordement réel, Litmus Edge pourra publier du JSON dans cette
arborescence ou utiliser Sparkplug B. Dans le second cas, un nœud de décodage
protobuf sera ajouté avant la validation ; le stockage et les tableaux de bord
ne changeront pas.

\section{Modèle de données}

\begin{figure}[H]
  \centering
  \includegraphics[width=0.93\textwidth]{modele-donnees.pdf}
  \caption{Organisation des données dans le bucket InfluxDB}
  \label{fig:model}
\end{figure}

La mesure \texttt{printer\_telemetry} regroupe les valeurs continues sous forme
de champs : niveau d'encre, niveau de solvant, viscosité, températures,
pression et compteurs. \texttt{printer\_event} conserve les changements d'état,
les défauts et les remplissages. \texttt{marking\_event} contient le DMC, le
résultat du contrôle et le grade de lecture. Les tags \texttt{printer\_id},
\texttt{line\_id}, \texttt{event\_type} et \texttt{verify\_grade} permettent les
filtres courants sans transformer chaque valeur numérique en série indexée.
Node-RED ajoute aussi le tag \texttt{classification} et les champs dérivés
\texttt{hours\_to\_empty}, \texttt{action\_days} et
\texttt{expiry\_limited} pour chaque fluide.

\section{Disponibilité liée à la demande}

\begin{figure}[H]
  \centering
  \includegraphics[width=0.92\textwidth]{logique-disponibilite.pdf}
  \caption{Distinction entre production, arrêt et attente normale}
  \label{fig:demand}
\end{figure}

La disponibilité est calculée sur les échantillons où la ligne demande une
impression :

\begin{equation}
A = \frac{N_{\text{demande et marche}}}
         {N_{\text{demande}}}\times 100
\end{equation}

Une période sans commande n'entre ni au numérateur ni au dénominateur. Si le
flux réel est publié à intervalle fixe, ce comptage approche correctement le
temps. Si Litmus publie uniquement lors d'un changement, il faudra reconstruire
des intervalles et pondérer par leur durée.

\section{Prévision et péremption}

Node-RED estime la consommation à partir de deux niveaux successifs. Pour
limiter l'effet du bruit, le débit observé est lissé avec le débit précédent :

\begin{equation}
r_k = 0{,}75\,r_{k-1} + 0{,}25\,
      \frac{L_{k-1}-L_k}{t_k-t_{k-1}}
\end{equation}

Pour un débit $r_k$ exprimé en points de pourcentage par heure et un niveau
courant $L_k$, le temps avant épuisement vaut :

\begin{equation}
T_{\text{vide}} = \frac{L_k}{r_k}
\end{equation}

La date d'action n'est pas forcément déterminée par ce calcul. Un fluide ouvert
peut expirer avant d'être vide. La date effective est :

\begin{equation}
D_{\text{effective}} =
\min\left(D_{\text{imprimée}},
D_{\text{ouverture}} + N_{\text{jours après ouverture}}\right)
\end{equation}

Une hausse supérieure à cinq points indique un remplissage et réinitialise
l'estimation. Le tableau de maintenance affiche les deux échéances et retient
celle qui arrive en premier.
